% Part: first-order-logic 
% Chapter: axiomatic-deduction 
% Section: rules-and-proofs

\documentclass[../../../include/open-logic-section]{subfiles}

\begin{document}

\iftag{FOL}
      {\olfileid{fol}{axd}{rul}}
      {\olfileid{pl}{axd}{rul}}

\olsection{القواعد و\usetoken{P}{derivation}}

\begin{explain}
  لعل أنساق ال!!{derivation}s البديهية أبسط أنساق ال!!{derivation}
  المنطقية. فال!!^a{derivation} متتالية من ال!!{formula}s لا غير.
  وإنما تعدّ !!a{derivation} إذا كانت كل !!{formula} فيها إما حالةً
  من بديهية، وإما مأخوذة بقاعدة استدلال من !!{formula}s سبقتها،
  واحدةً كانت أو أكثر. وال!!^a{derivation} !!{derive}s
  ال!!{formula} التي تختم بها المتتالية.
\end{explain}

\begin{defn}[!!^{derivability}]
إذا كانت~${\OLId{greek-variable}{Gamma}}$ مجموعة من !!{formula}s اللغة~$\Lang L$،
فالـ\emph{!!{derivation}} من~${\OLId{greek-variable}{Gamma}}$ متتالية منتهية~$!A_{\OLMathNumeral{1}}$،
\dots،~$!A_{\OLId{latin-ordinary}{n}}$ من ال!!{formula}s، لا يخلو فيها كل~${\OLId{latin-ordinary}{i}} \le {\OLId{latin-ordinary}{n}}$
من إحدى الحالات الآتية:
\begin{enumerate}
\item $!A_{\OLId{latin-ordinary}{i}} \in {\OLId{greek-variable}{Gamma}}$؛ أو
\item $!A_{\OLId{latin-ordinary}{i}}$ بديهية؛ أو
\item تلزم $!A_{\OLId{latin-ordinary}{i}}$ بقاعدة استدلال عن صيغة سابقة $!A_{\OLId{latin-ordinary}{j}}$، أو عن صيغتين
  سابقتين $!A_{\OLId{latin-ordinary}{j}}$ و$!A_{\OLId{latin-ordinary}{k}}$، على أن يكون ${\OLId{latin-ordinary}{j}}<{\OLId{latin-ordinary}{i}}$، وأن يكون ${\OLId{latin-ordinary}{k}}<{\OLId{latin-ordinary}{i}}$ في
  الحالة الثانية.
\end{enumerate}
\end{defn}

% تصحيح مصدر (0118-I1): عُرّف الدعم المنتهي الذي تُقاس به شروط الفروض المستعملة.
ونسمي ما أُخذ في ال!!{derivation} من~${\OLId{greek-variable}{Gamma}}$ بمقتضى الوجه الأول
\emph{دعم ال!!{derivation}}. وهو متناهٍ وإن لم تكن~${\OLId{greek-variable}{Gamma}}$ متناهية؛
فكل شرط يُردّ إلى مقدمات البرهان إنما يُردّ إلى هذا الدعم، لا إلى
ما لم يُستعمل من~${\OLId{greek-variable}{Gamma}}$.

وصحة ال!!{derivation} تتبع ما نجيزه من قواعد الاستدلال وما نختاره
من البديهيات. وقاعدة الاستدلال قضية شرطية تبيّن متى يجوز عدّ
الخطوة~$!A_{\OLId{latin-ordinary}{i}}$ في !!a{derivation} استدلالًا صحيحًا.


\begin{defn}[قاعدة الاستدلال]
\emph{قاعدة الاستدلال} بيان شرط كافٍ لصحة خطوة استدلال
في !!a{derivation} من~${\OLId{greek-variable}{Gamma}}$.
\end{defn}

فالمتتالية ذات العنصر الواحد~$!A$، متى كان~$!A \in {\OLId{greek-variable}{Gamma}}$،
تعدّ !!a{derivation} بنفسها. ومن ثم يصح أن نجعل الآتي
قاعدة استدلال أولى:
\begin{quote}
  إذا كان $!A \in {\OLId{greek-variable}{Gamma}}$، فإن $!A$ دائمًا خطوة استدلال صحيحة في كل
  !!{derivation} من~${\OLId{greek-variable}{Gamma}}$.
\end{quote}
وكذلك إذا كانت~$!A$ بديهية، كان~$!A$ وحده !!a{derivation}؛
فيصح أن نجعل الآتي قاعدة أخرى:
\begin{quote}
  إذا كانت $!A$ بديهية، فإن $!A$ خطوة استدلال صحيحة.
\end{quote}
وأما القواعد التي تنتفع بما سبق من !!{formula}s في المتتالية،
فمنها القاعدة الآتية، وتسمّى \emph{إثبات المقدَّم}:
\begin{quote}
  إذا ظهرت $!B \lif !A$ و$!B$ في موضع أعلى من ال!!{derivation}،
  فإن~$!A$ خطوة استدلال صحيحة.
\end{quote}
فإذا اقتصرنا على هذه القاعدة، صار تعريف ال!!{derivation}
المتقدّم إلى أن~$!A_{\OLMathNumeral{1}}$، \dots،~$!A_{\OLId{latin-ordinary}{n}}$ !!a{derivation}
إذا وفقط إذا تحقق لكل~${\OLId{latin-ordinary}{i}} \le {\OLId{latin-ordinary}{n}}$ واحد مما يأتي:
\begin{enumerate}
\item $!A_{\OLId{latin-ordinary}{i}} \in {\OLId{greek-variable}{Gamma}}$؛ أو
\item $!A_{\OLId{latin-ordinary}{i}}$ بديهية؛ أو
\item لبعض ${\OLId{latin-ordinary}{j}} < {\OLId{latin-ordinary}{i}}$ تكون $!A_{\OLId{latin-ordinary}{j}}$ هي $!B \lif !A_{\OLId{latin-ordinary}{i}}$، ولبعض ${\OLId{latin-ordinary}{k}} < {\OLId{latin-ordinary}{i}}$ تكون
  $!A_{\OLId{latin-ordinary}{k}}$ هي~$!B$.
\end{enumerate}
ومعنى الحالة الأخيرة أن~$!A_{\OLId{latin-ordinary}{i}}$ مأخوذة من~$!A_{\OLId{latin-ordinary}{j}}$
($!B \lif !A_{\OLId{latin-ordinary}{i}}$) و$!A_{\OLId{latin-ordinary}{k}}$ ($!B$) بإثبات المقدَّم.
فإذا فحصنا المتتالية من~${\OLMathNumeral{1}}$ إلى~${\OLId{latin-ordinary}{n}}$، ووجدنا في كل موضع
!!a{formula}~$!A_{\OLId{latin-ordinary}{i}}$ إما في~${\OLId{greek-variable}{Gamma}}$، وإما بديهية، وإما مما
تسوّغه قاعدة استدلال، صحّ عدّ المتتالية كلّها !!{derivation}.


\begin{defn}[!!^{derivability}]
ال!!{formula}~$!A$ \emph{!!{derivable}} من~${\OLId{greek-variable}{Gamma}}$،
وترميز ذلك~${\OLId{greek-variable}{Gamma}} \Proves !A$، متى وجد !!a{derivation}
من~${\OLId{greek-variable}{Gamma}}$ ينتهي إلى~$!A$.
\end{defn}

\begin{defn}[المبرهنات]
ال!!^a{formula}~$!A$ \emph{مبرهنة} متى وجد !!a{derivation}
لـ~$!A$ من المجموعة الخالية. فنكتب~$\Proves !A$ إذا كانت~$!A$
مبرهنة، ونكتب~$\Proves/ !A$ إذا لم تكن.
\end{defn}


\end{document}
